% Movimento Dependente: 2 blocos e 2 roldanas 
% Faro, sáb 06 dez 2025 20:54:15 
% Written by Tordar 
% pstricks2pdf.sh 2blocos_2roldanas
\documentclass{article}
\pagestyle{empty}

\usepackage{pst-all} 
\usepackage{pst-slpe}
\usepackage{graphicx}
\input{apnum}

% Roldana
\newcommand{\roldana}[5]{% x, y, radius, hole-radius, color
 \pscircle[linewidth=0.5mm,linecolor=lightgray,dimen=middle,fillstyle=radslope,slopebegin=gray,slopeend=white,slopesteps=50](#1,#2){#3}
  \ifnum\pdfstrcmp{#4}{0}=0\relax\else
    \pscircle[fillcolor=white,fillstyle=solid,linewidth=0.8pt](#1,#2){#4}
  \fi
  \pscircle[linewidth=1pt,linecolor=#5!50!black](#1,#2){#3}
}
% Bloco
\newcommand{\bloco}[5][]{% #1=options, #2=cx, #3=cy, #4=totalWidth, #5=height
  % Left frame (half of total width, full height)
  \psframe[linewidth=1pt,linecolor=gray,fillstyle=slope,slopebegin=white,slopeend=gray,slopesteps=15]
    (!#2 #4 2 div sub #3 #5 2 div sub)                     % lower-left of total area
    (!#2 #3 #5 2 div add)                                  % upper-right of left frame (center x, top)  
  % Right frame (half of total width, full height)
  \psframe[linewidth=1pt,linecolor=gray,fillstyle=slope,slopebegin=gray,slopeend=white,slopesteps=15]
    (!#2 #3 #5 2 div sub)                                  % lower-left of right frame (center x, bottom)
    (!#2 #4 2 div add #3 #5 2 div add)                     % upper-right of total area
  \psframe[#1](!#2 #4 2 div sub #3 #5 2 div sub)(!#2 #4 2 div add #3 #5 2 div add) 
}

\begin{document}

% Raio de cada roldana:
\evaldef\theri{0.75}
% Altura e largura de cada bloco: 
\evaldef\thew{1.5}
\evaldef\theh{1.5}
% Coordenadas do primeiro bloco: 
\evaldef\thexii{3}
\evaldef\theyii{3}
% Posição da primeira roldana: 
\evaldef\theyiii{7}
\evaldef\thexiii{\thexii+\theri}
% Posição da segunda roldana: 
\evaldef\theyv{5}
\evaldef\thexv{\thexii+3*\theri}
% Coordenadas do segundo bloco: 
\evaldef\xsb{\thexv}
\evaldef\ysb{1}
% Coordenadas auxiliares:
\evaldef\theyiv{\theyii+\theh/2}
\evaldef\xavi{\thexiii+\theri}
\evaldef\xavii{\thexv+\theri}
\evaldef\xaviii{\ysb+\theh/2}
\evaldef\xaix{\thexii-0.75*\thew}
\evaldef\xax{\xsb+0.75*\thew}

% Let's go:
\begin{pspicture}(0,0)(10,10)
% Teto: 
\psframe[linewidth=0pt,linecolor=lightgray,dimen=inner,fillstyle=hlines*,fillcolor=lightgray,hatchangle=45](1.5,9)(8.5,10) 
\psline[linewidth=2pt,linecolor=black](\thexiii,\theyiii)(\thexiii,9) 
% Blocos: 
\bloco[linewidth=2pt,linecolor=black]{\thexii}{\theyii}{\thew}{\theh}
\bloco[linewidth=2pt,linecolor=black]{\xsb}{\ysb}{\thew}{\theh}
% Roldanas e fios:
\roldana{\thexiii}{\theyiii}{\theri}{0.05}{gray} % 1 
\psline[linewidth=2pt,linecolor=black](\thexiii,\theyiii)(\thexiii,9)     % Fio teto      - roldana 1
\psline[linewidth=2pt,linecolor=black](\thexii,\theyiv)(\thexii,\theyiii) % Fio bloco   1 - roldana 1
\psline[linewidth=2pt,linecolor=black](\xavi,\theyiii)(\xavi,\theyv)      % Fio roldana 1 - roldana 2
\psline[linewidth=2pt,linecolor=black](\xavii,\theyv)(\xavii,9)           % Fio teto      - roldana 2
\roldana{\thexv}{\theyv}{\theri}{0.05}{gray} % 2 
\psline[linewidth=2pt,linecolor=black](\xsb,\theyv)(\xsb,\xaviii)         % Fio bloco 2   - roldana 2
\rput{0}(\xaix,\theyii){\scalebox{1.5}{$A$}}
\rput{0}(\xax,\ysb){\scalebox{1.5}{$B$}}
\end{pspicture}

\end{document}
